\section{Introduction}\label{sec:intro}
Hydration and food diaries require people to repeatedly notice an event, identify a quantity, and enter a record. Dietary self-monitoring has been studied extensively, but heterogeneous methods and adherence measures complicate interpretation~\cite{burke2011}. Lower interaction effort is therefore a design objective rather than evidence that a particular interface improves health. A laboratory or office offers an opportunity to instrument recurring routines without requiring a dedicated bottle to travel with every user.

Two problems remain when a scale replaces manual quantity entry. First, a decrease in container mass may reflect drinking, discarding water, spilling, or replacing the container. Commercial bottle evaluation has explicitly identified this distinction and the dependence on returning a bottle to a surface~\cite{cohen2022}. Second, a shared scale measures an object, not an authenticated person. Even accurate mass estimates can produce misleading personal histories if a user switches during an event. Conversational access adds another failure path: a fluent summary can conceal missing records, inconsistent units, or unauthorized queries.

\sys{} addresses these concerns as a proposed fixed-workplace research system. One scale acts as a drinking-cup station and the other supports food weighing. A Raspberry Pi 4 runs sensor acquisition; FastAPI mediates ingestion, authorization, event processing, and deterministic calculations; Supabase provides PostgreSQL storage. Discord supplies selection and correction interactions. An n8n agent chooses read-only analytics tools and explains returned evidence. We do not propose a new language model or a new load-cell transducer.

The study asks: \textbf{RQ1}, can temporal sensing and event confirmation improve hydration recording in the presence of ambiguous mass changes? \textbf{RQ2}, can versioned shared-device attribution and semi-automatic entry support reliable records with less interaction than manual entry? \textbf{RQ3}, can tool-grounded summaries preserve numeric correctness, authorization, and provenance?

The paper contributes (1) a sensing-to-record architecture that preserves uncertainty and correction history, (2) a method combining mass-based candidate events, food portion semantics, and session-based attribution, and (3) a prospective E1--E5 evaluation linking sensing and conversational failures to observable metrics. Tool use is motivated by ReAct and Toolformer~\cite{yao2023react,schick2023}; it does not by itself establish truthful or authorized behavior. All components below describe a proposed design. No formal E1--E5 measurements or independently verified running deployment are reported.
